\subsection{谱与预解算子}\label{sec:spectrum-and-resolvent}

本节将线性代数中的特征值问题推广到 Banach 空间上的算子. 主要目标之一是推广以下定理.

\begin{Theorem}
\label{thm:symmetric-matrix-spectral-basis}
设 $A$ 为 $n\times n$ 对称矩阵. 那么 $A$ 有 $n$ 个特征值 $\lambda_1,\ldots,\lambda_n$ (按代数重数计数), 且这些特征值全为实数. 此外, $\mathbb{R}^n$ 有一组标准正交基 $\{\bm{e}_1,\ldots,\bm{e}_n\}$, 使 $\bm{e}_i$ 为对应于 $\lambda_i$ 的特征向量.
\end{Theorem}

任何一本好的初等线性代数教材都给出了这一结果的证明. 这一结果将成为下面要证明的关于自伴紧算子的定理的推论.

首先要做的工作之一是推广特征值和特征向量的概念, 使之适用于本节所考察的算子 (这些算子可能定义在无穷维空间上, 也可能是无界的).

\begin{Definition}
\label{def:shifted-operator-and-resolvent}
设 $X$ 为复 Banach 空间, $(\mathcal{D}(A),A)$ 为从 $X$ 到 $X$ 的算子. 对任意 $\lambda\in\mathbb{C}$, 用下式定义算子 $(\mathcal{D}(A),A_\lambda)$:
\begin{equation}\label{eq:shifted-operator-definition}
A_\lambda:=A-\lambda I,
\end{equation}
其中 $I$ 为 $X$ 上的恒等算子. 若 $A_\lambda$ 有逆 (即它为单射), 则将其逆记为 $R_\lambda(A)$, 并称之为 $A$ 的\textbf{预解算子}.
\end{Definition}

\begin{Definition}
\label{def:resolvent-spectrum-and-spectral-parts}
设 $X\ne\{0\}$ 为复 Banach 空间, $(\mathcal{D}(A),A)$ 为从 $X$ 到 $X$ 的线性算子. 考察以下三个条件:\begin{enumerate}\item $R_\lambda(A)$ 存在,\item $R_\lambda(A)$ 有界,\item $R_\lambda(A)$ 的定义域在 $X$ 中稠密.\end{enumerate}将复平面 $\mathbb{C}$ 分解为以下两个集合.\begin{itemize}\item 算子 $A$ 的\textbf{预解集}为集合
\begin{equation}\label{eq:resolvent-set-definition}
\rho(A):=\{\lambda\in\mathbb{C}\mid\text{(1), (2), and (3) hold}\}.
\end{equation}
预解集中的元素 $\lambda\in\rho(A)$ 称为算子 $A$ 的\textbf{正则值}.\item 算子 $A$ 的\textbf{谱}是预解集的补集
\begin{equation}\label{eq:operator-spectrum-definition}
\sigma(A):=\mathbb{C}\backslash\rho(A).
\end{equation}
谱还可分解为三个互不相交的集合.\begin{itemize}\item \textbf{点谱}或\textbf{离散谱}为集合
\begin{equation}\label{eq:operator-point-spectrum-definition}
\sigma_p(A):=\{\lambda\in\sigma(A)\mid\text{(1) does not hold}\}.
\end{equation}
即点谱为使 $\mathcal{N}(A_\lambda)$ 非平凡的 $\lambda\in\mathbb{C}$ 所组成的集合. 点谱的元素称为\textbf{特征值}. 若 $\lambda\in\sigma_p(A)$, 则元素 $x\in\mathcal{N}(A_\lambda)$ 称为 $A$ 的\textbf{特征向量}或\textbf{特征函数}. $\mathcal{N}(A_\lambda)$ 的维数称为 $\lambda$ 的 (几何) \textbf{重数}.\item \textbf{连续谱}为集合
\begin{equation}\label{eq:operator-continuous-spectrum-definition}
\sigma_c(A):=\{\lambda\in\sigma(A)\mid\text{(1) and (3) hold but (2) does not}\}.
\end{equation}
\item \textbf{剩余谱}或\textbf{压缩谱}为集合
\begin{equation}\label{eq:operator-residual-spectrum-definition}
\sigma_r(A):=\{\lambda\in\sigma(A)\mid\text{(1) holds but (3) does not}\}.
\end{equation}
由于 $\overline{\mathcal{R}(A_\lambda)}\ne X$, 称值域被压缩了.\end{itemize}\end{itemize}
\end{Definition}

\begin{Definition}
\label{def:spectral-deficiency}
若 $X$ 为 Hilbert 空间, 则将 $\mathcal{R}(A_\lambda)^\perp$ 的维数称为 $\lambda\in\mathbb{C}$ 的\textbf{亏数}.
\end{Definition}

注意, 按这里的定义, $\lambda\in\sigma(A)$ 可以具有非零亏数而不属于压缩谱. 有些作者将压缩谱定义为所有亏数非零的 $\lambda\in\mathbb{C}$ 所组成的集合, 但这样一来, 点谱和压缩谱就未必互不相交.

\begin{Example}
\label{ex:finite-dimensional-spectrum-regular-eigenvalues}
线性代数的一个基本结果是, 对有限维空间上的线性算子 $A$, 其连续谱和压缩谱都是空集; 即复平面可分解为该算子的正则值和特征值.
\end{Example}

\begin{Example}
\label{ex:right-shift-compression-spectrum}
作为谱值并非特征值的算子的一个简单例子, 考察右移算子 $S_r:\ell^2\to\ell^2$. 复数 $\lambda=0$ 是谱的元素. 为看出这一点, 回忆预解算子 $R_0(S_r)$ 就是作用在定义域 $\{1,0,0,\ldots\}^\perp$ 上的左移算子 $S_l$. 虽然这个算子有界, 但其定义域在 $\ell^2$ 中不稠密. 因而 $\lambda=0$ 属于 $S_r$ 的压缩谱, 亏数为 $1$.
\end{Example}

谱理论是一个范围非常广且已有充分研究的课题. 这里的论述必然十分简略; 主要目标是建立特征函数展开这一工具. 因而先给出一个关于特征向量的基本定理.

\begin{Theorem}
\label{thm:distinct-eigenvalue-eigenvectors-independence}
若 $\lambda_i$, $i=1,\ldots,n$, 是算子 $(\mathcal{D}(A),A)$ 的互不相同的特征值, $x_i\in\mathcal{N}(A_{\lambda_i})$ 为对应的特征向量, 则集合\[\{x_1,x_2,\ldots,x_n\}\]线性无关.
\end{Theorem}

\begingroup
\postdisplaypenalty=10000
\renewcommand{\endProof}{\par\nobreak\hfill$\blacksquare$\par\vspace{0.1cm}}
\begin{Proof}
反设不然. 则存在整数 $k\in[2,n]$, 使集合 $\{x_1,\ldots,x_{k-1}\}$ 线性无关, 而 $x_k$ 可以用该集合展开; 即
\begin{equation}\label{eq:eigenvector-first-dependence-expansion}
x_k=\alpha_1x_1+\alpha_2x_2+\cdots+\alpha_{k-1}x_{k-1},
\end{equation}
其中系数 $\alpha_i$ 不全为零. 现在对等式两边作用 $(A-\lambda_kI)$, 得到\[\begin{aligned}
0&=(A-\lambda_kI)x_k\\
&=(A-\lambda_kI)[\alpha_1x_1+\alpha_2x_2+\cdots+\alpha_{k-1}x_{k-1}]\\
&=\alpha_1(\lambda_1-\lambda_k)x_1+\alpha_2(\lambda_2-\lambda_k)x_2\\
&\quad+\cdots+\alpha_{k-1}(\lambda_{k-1}-\lambda_k)x_{k-1}.
\end{aligned}\]由于 $\{x_1,\ldots,x_{k-1}\}$ 线性无关, 有
\begin{equation}\label{eq:eigenvector-independence-coefficient-identity}
(\lambda_i-\lambda_k)\alpha_i=0,\quad i=1,\ldots,k-1.
\end{equation}
然而, 由于 $\lambda_i\ne\lambda_m$, 这说明 $\alpha_i=0$, $i=1,\ldots,k-1$. 这产生矛盾, 证明完毕.
\end{Proof}
\endgroup
